\section{Przewidywany temat pracy}
Przewidywany temat pracy dyplomowej inżynierskiej to: \textbf{,,System sterowania bezzałogowego pojazdu podwodnego''}.

Praca ma charakter praktyczny, skupiający się na warstwie oprogramowania bezzałogowego pojazdu podwodnego (UUV -- \emph{Unmanned Underwater Vehicle}). Projektowany system ma w pierwszej fazie rozwoju obsługiwać pojazd uwiązany kablem (ROV -- \emph{Remotely Operated Vehicle}), jednak jego architektura logiczna i algorytmiczna jest rozwijana z myślą o pełnej skalowalności w stronę pojazdu autonomicznego (AUV -- \emph{Autonomous Underwater Vehicle}). 

Zakres inżynierski obejmuje sformułowanie nieliniowego modelu matematycznego w sześciu stopniach swobody, implementację i strojenie struktur sterowania oraz weryfikację działania w środowisku symulacyjnym.

\section{Cel pracy i jego uzasadnienie}

Celem podstawowym pracy dyplomowej jest opracowanie, implementacja oraz weryfikacja symulacyjna kompletnego systemu sterowania ruchem bezzałogowego pojazdu podwodnego, opartego na modelu matematycznym obiektu i przygotowanego do bezproblemowego wdrożenia na fizycznej platformie sprzętowej (\emph{sim-to-real transfer}).

\subsection{Cele szczegółowe}
\begin{enumerate}
    \item Opracowanie pełnego, nieliniowego modelu kinematyki i dynamiki pojazdu podwodnego z uwzględnieniem specyfiki środowiska morskiego (siły hydrodynamiczne, masa dodana, tłumienie wiskotyczne).
    \item Implementacja regulatorów, np. PID, wraz z algorytmem sterowania uwzględniającym specyficzną dynamikę drona.
    \item Utworzenie środowiska testowego w ROS 2 (\emph{Robot Operating System}) oraz symulatora Gazebo, umożliwiającego testy w pętli zamkniętej (\emph{Software-in-the-loop}).
    \item Wykorzystanie danych z układu IMU oraz czujnika ciśnienia (pomiar pośredni głębokości) do domknięcia pętli regulacji.
\end{enumerate}

\subsection{Uzasadnienie}
Projektowanie, budowa i weryfikacja bezzalogowych pojazdów podwodnych napotykają znacznie większe bariery implementacyjne niż ma to miejsce w przypadku bezzałogowych statków powietrznych. Środowisko wodne, ze względu na dużą gęstość, wprowadza silnie nieliniowe efekty dynamiczne oraz nieliniowe tłumienie hydrodynamiczne. 

Ze względu na brak dostępności sygnału GPS pod wodą, samolokalizacja opiera się na czujnikach inercyjnych (IMU) i hydrostatycznych. To sprawia, że mało jest dobrze udokumntowanych i dostępnych w domenie publicznej projektów wykorzystujących sterowanie inne niż manualne lub z wykorzystaniem prostych regulatorów.

\section{Przegląd stanu wiedzy}

Współczesne podejście do modelowania matematycznego jednostek pływających i podwodnych opiera się na standardach stowarzyszenia SNAME (\emph{Society of Naval Architects and Marine Engineers}), definiujących układy współrzędnych związanych z ziemią (NED) oraz z ciałem pojazdu. 

Kluczowym autorytetem w tej dziedzinie jest profesor Thor I. Fossen, którego macierzowy zapis równań ruchu uwzględnia sztywne ciało, siły bezwładności mas dodanych, macierze sil Coriolisa i sił odśrodkowych, a także wektory tłumienia wiskotycznego i sił hydrostatycznych (wypór i ciężar). 

Do weryfikacji poprawności obliczeń analitycznych powszechnie stosuje się otwarto-źródłowe pakiety symulacyjne, takie jak MSS (\emph{Marine Systems Simulator}) rozwijane m.in. w środowisku MATLAB.

\section{Metodyka}
Z uwagi na wysokie koszty budowy prototypów oraz trudną dostępność infrastruktury basenowej, kluczową rolę w nowoczesnej robotyce morskiej odgrywa faza symulacji typu SITL. Wykorzystanie ROS 2 w połączeniu z zaawansowanym fizycznie silnikiem Gazebo pozwala na generowanie danych syntetycznych oraz testowanie algorytmów sterowania w warunkach przypominających rzeczywiste. 

Odpowiednie odwzorowanie parametrów fizycznych w symulatorze umożliwia bezpośrednie przeniesienie wypracowanego kodu sterującego na mikrokontroler fizycznego robota bez ryzyka uszkodzenia sprzętu.

Choć praca koncentruje się na zagadnieniu oprogramowania, cała jej struktura jest projektowana pod kątem możliwości implementacji na rzeczywistym sprzęcie. Obliczenia, komunikacja oraz algorytmy sterowania będą realizowane na komputerze pomocniczym, np. Raspberry Pi.

\section{Plan pracy dyplomowej}
Harmonogram realizacji zadań inżynierskich obejmuje okres od kwietnia 2026 roku do stycznia 2027 roku i przedstawia się następująco:

\newpage

\begin{table}[hbtp]
    \centering
    \renewcommand{\arraystretch}{1.5} % Zwiększa odstępy wierszy o 50%
    \caption{Harmonogram realizacji zadań}
    \label{tab:harmonogram}
    \vspace{0.5em}
    \begin{tabular}{llp{10cm}}
        \toprule
        \textbf{Rok} & \textbf{Miesiąc} & \textbf{Zakres realizowanych zadań} \\
        \cmidrule(lr){1-1} \cmidrule(lr){2-2} \cmidrule(lr){3-3} % Podkreślenie tylko nagłówków
        2026 & Kwiecień - czerwiec   & Analiza i wykorzystanie literatury przedmiotu, szczegółowe studia nad stanem wiedzy w zakresie automatyki i inżynierii morskiej. \\
             & Lipiec     & Wyprowadzenie pełnego matematycznego modelu kinematyki i dynamiki pojazdu podwodnego w ujęciu 6 stopni swobody. \\
             & Sierpień   & Zbudowanie wirtualnego modelu robota oraz integracja środowiska ROS 2 + Gazebo. \\
             & Wrzesień   & Projektowanie i implementacja struktur sterowania: klasycznych struktur PID. \\
             & Październik & Strojenie parametrów układu regulacji, analiza uchybów ustalonych. \\
             & Listopad   & Przeprowadzenie kompleksowych testów symulacyjnych SITL, weryfikacja kodu. \\
             & Grudzień   & Opracowanie i analiza porównawcza wyników, redakcja tekstu pracy. \\
        \midrule % Delikatna linia przed ostatnim wierszem
        2027 & Styczeń    & Ostateczna redakcja tekstu pracy, naniesienie poprawek oraz złożenie gotowej pracy. \\
        \bottomrule
    \end{tabular}
\end{table}

\section{Podsumowanie}
Proponowany temat pracy odpowiada na istotne wyzwania współczesnej robotyki podwodnej, jakimi są nieliniowość środowiska oraz ograniczenia sensoryczne. 

Opracowanie zaawansowanego matematycznego modelu 6-DOF oraz implementacja struktury sterowania kompensującej dynamikę płynu pozwoli na uzyskanie znacznie wyższej precyzji prowadzenia pojazdu niż w przypadku standardowych podejść liniowych. 

Oprogramowanie będzie w pełni gotowe do bezproblemowej migracji z symulatora na fizycznego drona, co stanowi solidny punkt wyjścia do budowy autonomicznych systemów klasy AUV.